\section{Examples of partial differential equations}

A partial differential equation relates an unknown function of several independent variables to its partial derivatives. The independent variables often include position $\mathbf{x}$ and time $t$, but they may instead be spatial coordinates, momenta, or parameters. Some PDEs describe a single scalar field $u$, whereas others are systems for vector- or tensor-valued fields. Equations without a time variable are called stationary; equations containing time derivatives describe evolution.

The same equation can occur in very different settings. Conversely, changing physical assumptions may lead to a different PDE for the same broad phenomenon. The examples below provide a working catalogue before the detailed derivations in the rest of this chapter.

\subsection{Transport and conservation}

The first-order transport equation
\begin{equation}
u_t+c u_x=0
\end{equation}
describes a profile transported at constant speed $c$. Its solutions have the form $u(x,t)=g(x-ct)$, so the profile moves without changing shape. It can model the concentration of a tracer carried by a uniform one-dimensional flow.

More generally, a conserved density $\rho(\mathbf{x},t)$ with flux $\mathbf{J}$ satisfies the continuity equation
\begin{equation}
\rho_t+\nabla\cdot\mathbf{J}=s,
\end{equation}
where $s$ is a source or sink. For a fluid with velocity $\mathbf{v}$, $\mathbf{J}=\rho\mathbf{v}$, and conservation of mass gives
\begin{equation}
\rho_t+\nabla\cdot(\rho\mathbf{v})=0.
\end{equation}
The divergence term measures net outward flux from a small volume. This balance-law viewpoint also underlies heat conduction and charge conservation.

\subsection{Waves and elasticity}

The scalar wave equation in $d$ spatial dimensions is
\begin{equation}
u_{tt}-c^2\nabla^2u=0.
\end{equation}
It models small disturbances that propagate at speed $c$, including sound in a uniform medium and appropriate components of elastic motion. The one-dimensional version is the vibrating-string equation derived below.

For a homogeneous isotropic elastic solid, the displacement $\mathbf{u}(\mathbf{x},t)$ is vector valued. A standard linearized model is
\begin{equation}
\rho\mathbf{u}_{tt}
=(\lambda+\mu)\nabla(\nabla\cdot\mathbf{u})+\mu\nabla^2\mathbf{u},
\end{equation}
where $\lambda$ and $\mu$ are Lam\'e parameters. The two terms describe different deformation modes: compressional motion is associated with changes of volume, while shear motion is associated with changes of shape.

Flexural motion requires higher spatial derivatives. An ideal slender beam may satisfy
\begin{equation}
u_{tt}+K u_{xxxx}=0,
\end{equation}
and an ideal thin plate may satisfy
\begin{equation}
u_{tt}+K\nabla^4u=0,
\qquad \nabla^4=(\nabla^2)^2.
\end{equation}
Unlike a taut string, a beam or plate resists bending, which is why the bending equations contain fourth-order derivatives.

\subsection{Diffusion, dispersion, and quantum evolution}

The heat equation
\begin{equation}
u_t=\kappa\nabla^2u
\end{equation}
describes the smoothing of a temperature or concentration field. The same equation appears in diffusion models, although the variable $u$ then represents concentration or probability density rather than temperature. A convection--diffusion equation adds directed transport:
\begin{equation}
u_t+\mathbf{v}\cdot\nabla u=\kappa\nabla^2u.
\end{equation}

The time-dependent Schr\"odinger equation has a superficially similar spatial operator but a fundamentally different interpretation:
\begin{equation}
i\hbar\psi_t=
\left(-\frac{\hbar^2}{2m}\nabla^2+V(\mathbf{x})\right)\psi.
\end{equation}
Here $\psi$ is a complex wavefunction, $m$ is mass, and $V$ is a potential energy. Diffusion decreases gradients and is irreversible in its usual physical direction of time; Schr\"odinger evolution preserves quantum probability and is reversible for a time-independent Hamiltonian.

For a time-independent potential, solutions with a definite energy have the form $\psi(\mathbf{x},t)=\phi(\mathbf{x})e^{-iEt/\hbar}$. Substitution gives the stationary Schr\"odinger eigenvalue problem
\begin{equation}
\left(-\frac{\hbar^2}{2m}\nabla^2+V(\mathbf{x})\right)\phi=E\phi.
\end{equation}
Boundary conditions and the requirement that $\phi$ be normalizable select the permitted energy levels $E$. Thus a time-dependent PDE becomes a stationary spectral problem for the spatial wavefunction $\phi$.

Relativistic quantum mechanics uses the Dirac equation for a four-component complex spinor $\Psi$:
\begin{equation}
\left(i\hbar\gamma^\mu\partial_\mu-mc\right)\Psi=0.
\end{equation}
The matrices $\gamma^\mu$ encode the coupling between spin and spacetime. This is a first-order linear system in space and time; applying a conjugate Dirac operator shows its connection with the second-order relativistic wave equation.

Dispersive waves can be represented, in one simple linear model, by
\begin{equation}
u_t+\gamma u_{xxx}=0.
\end{equation}
Different wavelengths travel at different speeds, so an initially localized pulse can change shape even in a linear medium. This distinguishes dispersion from the nondispersive transport equation.

\subsection{Electromagnetic fields}

In vacuum, Maxwell's equations form a coupled system for electric and magnetic fields:
\begin{align}
\nabla\cdot\mathbf{E}&=\frac{\rho_e}{\varepsilon_0},
&\nabla\cdot\mathbf{B}&=0,\\
\nabla\times\mathbf{E}&=-\mathbf{B}_t,
&\nabla\times\mathbf{B}&=\mu_0\mathbf{J}+\mu_0\varepsilon_0\mathbf{E}_t.
\end{align}
The divergence equations constrain the allowed fields, while the curl equations describe their time evolution. In a charge-free, current-free region, taking a curl and using the vector identity from Chapter 2 shows that both fields satisfy wave equations with speed
\begin{equation}
c=\frac{1}{\sqrt{\mu_0\varepsilon_0}}.
\end{equation}

Yang--Mills theory generalizes Maxwell theory to a gauge potential $A_\mu$ whose components take values in a noncommutative Lie algebra. With covariant derivative $D_\mu=\partial_\mu+gA_\mu$, the curvature is
\begin{equation}
F_{\mu\nu}=\partial_\mu A_\nu-\partial_\nu A_\mu
+g[A_\mu,A_\nu],
\end{equation}
and the source-free Yang--Mills equation is
\begin{equation}
D_\mu F^{\mu\nu}=0.
\end{equation}
This nonlinear system reduces formally to the vacuum Maxwell equations when the field components commute. It supplies the classical field equations for non-Abelian gauge theories.

\subsection{Fluid flow and nonlinear equations}

For an incompressible viscous fluid of constant density $\rho$, velocity $\mathbf{v}$, pressure $p$, and dynamic viscosity $\mu$, the Navier--Stokes equations are
\begin{equation}
\rho\left(\mathbf{v}_t+(\mathbf{v}\cdot\nabla)\mathbf{v}\right)
=-\nabla p+\mu\nabla^2\mathbf{v},
\qquad \nabla\cdot\mathbf{v}=0.
\end{equation}
The nonlinear convection term $(\mathbf{v}\cdot\nabla)\mathbf{v}$ is the acceleration experienced while moving with the fluid. The pressure enforces incompressibility rather than following an independent scalar evolution equation. Neglecting viscosity produces the Euler equations for an ideal fluid.

A simpler nonlinear conservation law is the inviscid Burgers equation
\begin{equation}
u_t+u u_x=0.
\end{equation}
Its characteristic speed is the solution value $u$ itself. Faster portions of a profile can catch slower portions, making a smooth classical solution fail and motivating weak solutions and shock conditions.

\subsection{Spacetime geometry}

General relativity describes gravity through a Lorentzian spacetime metric $g_{\mu\nu}$. Einstein's field equation is the nonlinear tensor system
\begin{equation}
G_{\mu\nu}+\Lambda g_{\mu\nu}
=\frac{8\pi G}{c^4}T_{\mu\nu},
\end{equation}
where $G_{\mu\nu}=R_{\mu\nu}-\tfrac12 Rg_{\mu\nu}$ is the Einstein tensor, $T_{\mu\nu}$ is the stress--energy tensor, $G$ is Newton's gravitational constant, and $\Lambda$ is the cosmological constant. The Ricci tensor $R_{\mu\nu}$ contains second derivatives of $g_{\mu\nu}$ as well as terms quadratic in its first derivatives.

Because the metric also determines the coordinate geometry, the equations possess coordinate freedom rather than being an ordinary system of independent scalar PDEs. After a coordinate or gauge condition is imposed, their evolution form is hyperbolic in suitable circumstances. Vacuum gravitational waves arise from the source-free equation $G_{\mu\nu}=0$ and are small propagating perturbations of the metric.

\subsection{A diffusion equation outside physics}

The Black--Scholes equation illustrates that PDEs also arise in mathematical finance. Under its idealized assumptions, the value $V(S,t)$ of a European derivative satisfies
\begin{equation}
V_t+\frac12\sigma^2S^2V_{SS}+rS V_S-rV=0,
\end{equation}
where $S$ is the underlying asset price, $r$ is a risk-free interest rate, and $\sigma$ is volatility. Mathematically, this is a parabolic equation with variable coefficients. Its terminal payoff plays a role analogous to initial data, although the equation is commonly solved backward from the expiration time.

These examples illustrate why classification, initial data, boundary data, and coordinate choice are indispensable. The following sections develop the wave, heat, Laplace, and Poisson equations in detail.